\section{Exact current-to-link and plaquette
audit}\label{fabel-current-to-lattice-audit-md-exact-current-to-link-and-plaquette-audit}

This companion continues the Fabel/Jacobi and Cartan gauge map with the
original lattice coefficients. It computes the actual configuration
observable induced by a regular fluid segment and records why a nonzero
sourced classical field is not a source-free Yang--Mills vacuum state.

\subsection{1. Source current retained in physical
units}\label{fabel-current-to-lattice-audit-md-source-current-retained-in-physical-units}

Use (X\^{}0=ct), (X\textsuperscript{i=x}i), (T=-i\sigma\emph{3/2),
(A\_0=0), (A\_i=\lambda u\_iT), and (D\^{}\mu F}\{\mu\nu\}=g\^{}2j\_\nu)
with metric (operatorname\{diag\}(-1,1,1,1)). For an incompressible
fluid, \[
j_0=0,qquad
j_i=\frac{\lambda}{g^2}\left(\Delta u_i-c^{-2}\partial_t^2u_i\right)T.
\] Differentiating the full positive-viscosity equation once gives, with
(w\_i=\partial\_tu\_i), \[
j_i=\frac{\lambda}{g^2}\left[\Delta u_i-c^{-2}\left(\nu\Delta w_i-u_j\partial_jw_i-w_j\partial_ju_i-\partial_i\partial_tp+\partial_tf_i\right)\right]T.
\] Every coefficient and sign is retained. The current is covariantly
conserved on the Cartan image, but it is generally nonzero. The
stationary Fabel field (W) has (w=0), so its current reduces to
(j\_i=(\lambda/g\^{}2)\Delta W\_iT); the source record also gives the
nonzero viscous curl residual ((0,0,-18\nu t)) on the (x=y=0) axis. Thus
the polynomial escape is not a source-free positive-viscosity solution.

\subsection{2. Strong non-Abelian link
map}\label{fabel-current-to-lattice-audit-md-strong-non-abelian-link-map}

Let (mathcal A) be any smooth (\mathfrak{su}(2))-connection on a cubical
spatial mesh of spacing (a\textgreater0), with oriented edge (e=(n,i))
from (x\_n=an) to (x\_n+a e\_i). Define the exact parallel-transport
link \[
U_e[\mathcal A]=\mathcal P\exp\left(-\int_0^a\mathcal A_i(x_n+s e_i,t)\,ds\right)\in\mathrm{SU}(2).
\] For a smooth vertex gauge function (h), the ODE for parallel
transport gives the exact endpoint law \[
U_e[\mathcal A^h]=h(x_n)^{-1}U_e[\mathcal A]h(x_n+a e_i).
\] This is the required typed map from continuum connections modulo
gauge to lattice edge variables with the original orientation. It
retains all non-Abelian commutators through path ordering.

For the Cartan connection, path ordering is unnecessary and \[
U_i(n,t)=\exp\left(-\lambda T\int_0^a u_i(an+s e_i,t)\,ds\right).
\] The simpler sampled expression (exp(-a\lambda u\_i(an,t)T)) is only a
first-order approximation and is not used for exact gauge transport.

\subsection{3. Exact plaquette expansion and Hamiltonian
scale}\label{fabel-current-to-lattice-audit-md-exact-plaquette-expansion-and-hamiltonian-scale}

Let
(U\_p=U\_i(n)U\_j(n+e\_i)U\_i(n+e\_j)\textsuperscript{\{-1\}U\_j(n)}\{-1\})
for (i\textless j). For a (C\^{}2) connection, the non-Abelian transport
theorem gives \[
U_p=I-a^2F_{ij}(x_n,t)+O(a^3)
\] in operator norm on every compact region on which the (C\^{}1)-norm
of (mathcal A) and the (C\^{}0)-norm of (F) are bounded. The (O(a\^{}3))
constant is local and is not claimed uniform along a concentrating
sequence.

On the Cartan image, (F\_\{ij\}=\lambda(\partial\_i u\_j-\partial\_j
u\_i)T), and the plaquette is exactly a single Cartan exponential with
angle \[
\vartheta_p= -\lambda\oint_{\partial p}u\cdot dx.
\] Since (operatorname\{tr\}exp(\vartheta T)=2\cos(\vartheta/2)), the
original magnetic term is \[
b(2-W_p)=\frac1{2g^2a}\,\bigl(2-2\cos(\vartheta_p/2)\bigr)
 =\frac{\vartheta_p^2}{8g^2a}+O(\vartheta_p^4/a).
\] For a regular fixed field, (artheta\_p=-a\^{}2\lambda(\partial\_i
u\_j-\partial\_j u\_i)+O(a\^{}3)), hence one plaquette contributes \[
\frac{a^3\lambda^2}{8g^2}(\partial_i u_j-\partial_j u_i)^2+O(a^4)+O(a^7|F|^4).
\] Summing all three spatial orientations and taking a fixed smooth
compact region gives the Riemann limit \[
\sum_{p}b(2-W_p)\longrightarrow
\frac{\lambda^2}{8g^2}\int |\nabla\times u|^2\,d^3x,
\] with the displayed normalization (-2\operatorname{tr}(T\^{}2)=1).
This is a magnetic energy, not a mass eigenvalue.

The electric term remains \[
\kappa\sum_eE_e,qquad \kappa=\frac{2g^2}{a},
\] and depends on the wave-functional derivatives in the physical
Hilbert space. A classical time derivative in (j\_i) is not an electric
quantum excitation until a state and its operator domain are specified.

\subsection{4. Consequence for the Fabel
direction}\label{fabel-current-to-lattice-audit-md-consequence-for-the-fabel-direction}

The exact Fabel material map sends (e\_3) to
(zeta\_\tau=A\_\tau\^{}\{-\mathsf T\}e\_3) with
(\textbar{}\zeta\emph{\tau\textbar{}\sim(9/4)(1-8\tau)\^{}\{-3/2\}).
This is a tangent/cotangent amplification. The link map above sends an
actual connection (\mathcal A(u)) to (U\_e{[}\mathcal A{]}); there is no
typed identity sending (zeta}\tau) directly to (u), (F), or a
wave-functional derivative. If an independently proved fluid profile has
(\textbar{}\nabla\times u\textbar{}\to\infty) in a region whose lattice
cells resolve that region, the magnetic form grows according to the
preceding integral (or leaves the uniform Taylor regime). Either outcome
is a high-energy or loss-of-control signal. It does not yield a sequence
with excitation quotient tending to zero.

The source current has a second independent obstruction. The source-free
Hamiltonian (H=\kappa\sum\_eE\_e+b\sum\_p(2-W\_p)) contains no external
(j\_i). A sourced classical connection can be coupled only after
enlarging the theory and proving the corresponding Gauss constraint,
vacuum, and positive form. Dropping (j\_i) changes the equation;
retaining it does not produce a vector in the source-free vacuum Hilbert
space.

\subsection{5. Physical projection and raw gap
test}\label{fabel-current-to-lattice-audit-md-physical-projection-and-raw-gap-test}

For a lattice functional (v(U)), the physical vector is
(P\_\{\mathcal G\}(v\psi)) in (L\^{}2(\lambda)), or equivalently
(P\_\{\mathcal G\}v) in (L\textsuperscript{2(\psi}2\lambda)). The exact
projection identity is \[
\|v\|_{\psi^2\lambda}^2=\|P_{\mathcal G}v\|^2+\|(I-P_{\mathcal G})v\|^2.
\] The link map gives a configuration observable, but it does not prove
(P\_\{\mathcal G\}(v\psi)\ne0), vacuum orthogonality, form-domain
membership, or a small quotient. A mass-gap violation still requires
nonzero (v\_n\perp\psi) with \[
\frac{q_H[v_n]}{\|v_n\|^2_{\psi^2\lambda}}\to0,
\] where (q\_H) is the unchanged excitation form. Neither the sign
(a=-1/4) nor curvature concentration supplies this inequality.

The exact conclusion is therefore a complete continuation of the
classical field into gauge-covariant lattice observables, together with
a proved energy/current obstruction. The Fabel/Jacobi data still do not
produce the required physical vacuum-orthogonal low-energy sequence.
